\documentclass[14pt,a4paper]{scrreprt}

\usepackage{cmap}
\usepackage[T1]{fontenc} 
\usepackage[utf8]{inputenc}
\usepackage[english,russian]{babel}

\usepackage{caption}
\usepackage{subcaption}

\usepackage{float}

\usepackage{enumitem}

\usepackage{graphicx}
\usepackage{multirow}


\usepackage{pgfplots}
\pgfplotsset{compat=newest}
\usepgfplotslibrary{units}

\usepackage{longtable}

\usepackage{caption}
\captionsetup{labelsep=endash}
\captionsetup[figure]{name={Рисунок}}
\captionsetup[subtable]{labelformat=simple}
\captionsetup[subfigure]{labelformat=simple}
\renewcommand{\thesubtable}{\text{Таблица }\arabic{chapter}\text{.}\arabic{table}\text{.}\arabic{subtable}\text{ --}}
\renewcommand{\thesubfigure}{\text{Рисунок }\arabic{chapter}\text{.}\arabic{figure}\text{.}\arabic{subfigure}\text{ --}}


\usepackage{textcomp}

\usepackage{amsmath}
\usepackage{amsfonts}
\usepackage{array}

\usepackage{geometry}
\geometry{left=30mm}
\geometry{right=15mm}
\geometry{top=20mm}
\geometry{bottom=20mm}
\geometry{foot=1.7cm}

\usepackage{titlesec}
\titleformat{\section}
{\normalsize\bfseries}
{\thesection}
{1em}{}
\titlespacing*{\chapter}{0pt}{-30pt}{8pt}
\titlespacing*{\section}{\parindent}{*4}{*4}
\titlespacing*{\subsection}{\parindent}{*4}{*4}

% Маркировка для списков
\def\labelitemi{$\circ$}
\def\labelitemii{$*$}


\usepackage{setspace}
\onehalfspacing % Полуторный интервал

\frenchspacing
\usepackage{indentfirst} % Красная строка

\usepackage{titlesec}
\usepackage{xcolor}
% Названия глав
\titleformat{\section}{\normalsize\textmd}{\thesection}{1em}{}

\definecolor{gray35}{gray}{0.35}

\newcommand{\hsp}{\hspace{20pt}} % длина линии в 20pt

\titleformat{\chapter}[hang]{\Huge}{\textcolor{gray35}{\thechapter.}\hsp}{0pt}{\Huge\textmd}

\titleformat{\section}{\Large}{\textcolor{gray35}\thesection}{20pt}{\Large\textmd}
\titleformat{\subsection}{\Large}{\thesubsection}{20pt}{\Large\textmd}
\titleformat{\subsubsection}{\normalfont\textmd}{}{0pt}{}

% Настройки введения

\addtocontents{toc}{\setcounter{tocdepth}{2}}
\addtocontents{toc}{\setcounter{secnumdepth}{1}}

\usepackage{tocloft,lipsum,pgffor}

\addtocontents{toc}{~\hfill\textnormal{Страница}\par}

\renewcommand{\cftpartfont}{\normalfont\textmd}

\addto\captionsrussian{\renewcommand{\contentsname}{Содержание}}
\renewcommand{\cfttoctitlefont}{\Huge\textmd}

\renewcommand{\cftchapfont}{\normalfont\normalsize}
\renewcommand{\cftsecfont}{\normalfont\normalsize}
\renewcommand{\cftsubsecfont}{\normalfont\normalsize}
\renewcommand{\cftsubsubsecfont}{\normalfont\normalsize}

\renewcommand{\cftchapleader}{\cftdotfill{\cftdotsep}}

\usepackage{listings}
\usepackage{xcolor}

\usepackage[pdftex]{hyperref} % Гиперссылки
\hypersetup{hidelinks}

% Листинги 
\usepackage{listings}

\definecolor{darkgray}{gray}{0.15}

\definecolor{teal}{rgb}{0.25,0.88,0.73}
\definecolor{gray}{rgb}{0.5,0.5,0.5}
\definecolor{b-red}{rgb}{0.88,0.25,0.41}
\definecolor{royal-blue}{rgb}{0.25,0.41,0.88}

\usepackage{listings}
\lstset{
	aboveskip=3mm,
	belowskip=3mm,
	frame=tb,
	frame=single,
	basicstyle=\footnotesize\ttfamily,
	numberstyle=\tiny\color{gray},
	keywordstyle=\color{royal-blue},
	commentstyle=\color{gray35},
	stringstyle=\color{b-red},
	numbers=left,
	numbersep=5pt,
	numberstyle=\tiny,
	showstringspaces=false, 
	captionpos=t,
	tabsize=4,
	language=Lisp
}

% какой то сложный кусок со стак эксчейндж для квадратных скобок
\makeatletter
\newenvironment{sqcases}{%
	\matrix@check\sqcases\env@sqcases
}{%
	\endarray\right.%
}
\def\env@sqcases{%
	\let\@ifnextchar\new@ifnextchar
	\left\lbrack
	\def\arraystretch{1.2}%
	\array{@{}l@{\quad}l@{}}%
}
\makeatother

% и для матриц
\makeatletter
\renewcommand*\env@matrix[1][\arraystretch]{%
	\edef\arraystretch{#1}%
	\hskip -\arraycolsep
	\let\@ifnextchar\new@ifnextchar
	\array{*\c@MaxMatrixCols c}}
\makeatother


\begin{document}

\begin{titlepage}
	
	\newgeometry{left=2cm, right=2cm, top=2.5cm, bottom=2.5cm}
	\fontsize{12pt}{12pt}\selectfont
	
	\noindent \begin{minipage}{0.13\textwidth}
		\includegraphics[width=\linewidth]{assets/bmstu-logo.png}
	\end{minipage}
	\noindent\begin{minipage}{0.85\textwidth}\centering
		\textbf{\textsc{Министерство науки и высшего образования Российской Федерации}}\\
		\textbf{\textsc{Федеральное государственное бюджетное образовательное 	учреждение высшего образования}}\\
		\textbf{\textsc{Московский государственный технический университет имени 	Н.Э.~Баумана}}\\
		\textbf{\textsc{(национальный исследовательский университет)}}\\
		\textbf{\textsc{(МГТУ им. Н.Э. Баумана)}}\\
	\end{minipage}
	
	\noindent\rule{18cm}{1.5pt}
	
	\vspace{8mm}
	
	\noindent\textnormal{ФАКУЛЬТЕТ}\hspace{5mm} \underline{\textnormal{~~~~~~~~~~~~~~~~~~«Информатика и системы управления»~~~~~~~~~~~~~~~~~~}} \newline\newline
	\textnormal{КАФЕДРА}\hspace{5mm} \underline{\textnormal{~~«Программное обеспечение ЭВМ и информационные технологии»~~}}
	\newline\newline
	\textnormal{НАПРАВЛЕНИЕ ПОДГОТОВКИ}\hspace{5mm} \underline{\textnormal{~~~~~~«09.03.04 Программная инженерия»~~~~~~~~}}
	
	\vspace{2.5cm}
	
	\begin{center}
		\Large\textbf{\textsc{ОТЧЕТ}}\\
		\Large\textbf{\textsc{ПО ЛАБОРАТОРНОЙ РАБОТЕ №18}}\\
	\end{center}
	
	\vspace{1cm}
	
	\noindent\textnormal{Название:} \hspace{15mm} \underline{\textnormal{~~~~~~Формирование и модификация списков на Prolog~~~~~~~}}\noindent
	
	\vspace{1.3cm}
	
	\noindent\textnormal{Дисциплина:} \hspace{10mm} \underline{\textnormal{~~~~~~Функциональное и логическое программирование~~~~~~}}\noindent
	
	\vspace{1.5cm}
	
	\noindent\textnormal{Студент} \hspace{17mm}
	\underline{\textnormal{{~~~~ИУ7-64Б~~~}}}
	\hspace{20mm}
	\underline{\textnormal{\hphantom{~~~~~~~~~~~~~~~~~~~~~~~~~~~}}} \hspace{14mm}
	\underline{\textnormal{~~~С. Д. Параскун~~~}}
	
	\vspace{2mm}
	\noindent\textnormal{\hphantom{Студент}} \hspace{23mm}\noindent
	\fontsize{8pt}{8pt}
	\textnormal{Группа}\hspace{40mm}\textnormal{Подпись, дата} \hspace{30mm}\noindent\textnormal{И. О. Фамилия}
	
	\vspace{0.5cm}
	
	\fontsize{12pt}{12pt}\selectfont
	\noindent\textnormal{Преподаватель} \hspace{52mm}
	\underline{\textnormal{\hphantom{~~~~~~~~~~~~~~~~~~~~~~~~~~~}}} \hspace{14mm}
	\noindent\underline{\textnormal{~~Н. Б. Толпинская~~}}
	
	\vspace{2mm}
	\noindent\textnormal{\hphantom{Студент}} \hspace{17mm}\noindent
	\fontsize{8pt}{8pt}
	\hphantom{Группа}\hspace{43mm}\textnormal{Подпись, дата} \hspace{30mm}\noindent\textnormal{И. О. Фамилия}
	
	\vspace{0.5cm}
	
	\fontsize{12pt}{12pt}\selectfont
	\noindent\textnormal{Преподаватель} \hspace{52mm}
	\underline{\textnormal{\hphantom{~~~~~~~~~~~~~~~~~~~~~~~~~~~}}} \hspace{14mm}
	\noindent\underline{\textnormal{~~Ю. В. Строганов~~}}
	
	\vspace{2mm}
	\noindent\textnormal{\hphantom{Студент}} \hspace{17mm}\noindent
	\fontsize{8pt}{8pt}
	\hphantom{Группа}\hspace{43mm}\textnormal{Подпись, дата} \hspace{30mm}\noindent\textnormal{И. О. Фамилия}
	
	\vspace{1cm}
	
	\fontsize{12pt}{12pt}\selectfont
	
	\begin{center}
		\vfill
		Москва, ~\the\year
		~г.
	\end{center}
	\restoregeometry

\end{titlepage}


\thispagestyle{empty}

\chapter{Практическое задание}

\section{Задание}

Используя хвостовую рекурсию, разработать программу, позволяющую найти:
\begin{itemize}
	\item $n!$;
	\item $n$-е число Фибоначчи.
\end{itemize}
Убедиться в правильности результатов.

Для одного из вариантов вопроса и каждого задания составить таблицу, отражающую конкретный порядок работы системы:

\section{Код программы}

\begin{lstlisting}[language=Prolog]
predicates

	factorialIter(integer, integer, integer).
	factorial(integer, integer).
	fibonIter(integer, integer, integer, integer).
	fibon(integer, integer).

clauses

	factorialIter(N, Res, IterRes) :- N = 0, !, Res = IterRes.
	factorialIter(N, Res, IterRes) :- N > 0, NewN = N - 1, 
		NewIterRes = N * IterRes,
factorialIter(NewN, Res, NewIterRes). 
	factorial(N, Res) :- factorialIter(N, Res, 1).
	
	fibonIter(N, Res, FNum, _) :- N = 0, !, Res = FNum.
	fibonIter(N, Res, _, SNum) :- N = 1, !, Res = SNum.
	fibonIter(N, Res, FNum, SNum) :- N > 1, NewN = N - 1, 
		NewFNum = SNum, NewSNUm = FNum + SNum, 
	fibonIter(NewN, Res, NewFNum, NewSNum).
	fibon(N, Res) :- fibonIter(N, Res, 0, 1).

goal

	factorial(4, Res).
	%fibon(7, Res).
\end{lstlisting}

\section{Таблицы выполнения программы}

Запрос для задания factorial(3, Res).
\begin{table}[H]
	\begin{tabular}{|p{0.8cm\small}|p{4.7cm\small}|p{5.7cm\small}|p{4cm\small}|}	
		\hline
		№ шага & Состояния резольвенты и вывод: дальнейшие действия (почему?) & Для каких термов запускается алгоритм унификации: T1=T2 и каков результат (подстановка) & Дальнейшие действия: прямой ход или откат (почему и к чему приводит?)\\
		\hline
		1 & Резольвента:\begin{itemize} \item factorial(3, Res)  \end{itemize} & factorial(3, Res) = factorialIter(N, Res, IterRes). Унификация неуспешна & Прямой ход, переход к следующему предложению \\
		\hline
		2 & ... & ... & ...\\
		\hline
		3 & Резольвента:\begin{itemize} \item factorialIter(3, Res, 1) \end{itemize} & factorial(3, Res) = factorial(N, Res). Унификация успешна\linebreak \textbf{Подстановка:} \{N=3, Res=Res\} & Прямой ход, решение цели резольвенты \\
		\hline
		4 & Резольвента:\begin{itemize} \item  3 = 0, \item !, \item Res = 1, \item ! \end{itemize} & factorialIter(3, Res, 1) = factorialIter(N, Res, IterRes). Унификация успешна\linebreak \textbf{Подстановка:} \{N=3, Res=Res, IterRes=1\} & Прямой ход, решение цели резольвенты\\
		\hline
		5 & Резольвента:\begin{itemize} \item !, \item Res = 1, \item ! \end{itemize} & 3 = 0. Ложь & Откат, переход к следующему шагу относительно 4\\
		\hline
	\end{tabular}
\end{table}

\begin{table}[H]
	\begin{tabular}{|p{0.8cm\small}|p{4.7cm\small}|p{5.7cm\small}|p{4cm\small}|}	
		\hline
		№ шага & Состояния резольвенты и вывод: дальнейшие действия (почему?) & Для каких термов запускается алгоритм унификации: T1=T2 и каков результат (подстановка) & Дальнейшие действия: прямой ход или откат (почему и к чему приводит?)\\
		\hline
		6 & Резольвента:\begin{itemize} \item 3 > 0, \item NewN = 3 - 1, 
			\item NewIterRes = 3 * 1,
\item factorialIter(NewN, Res, NewIterRes) \end{itemize} & factorialIter(3, Res, 1) = factorialIter(N, Res, IterRes). Унификация успешна\linebreak \textbf{Подстановка:} \{N=3, Res=Res, IterRes=1\} & Прямой ход, решение цели резольвенты\\
		\hline
		7 & Резольвента:\begin{itemize} \item NewN = 3 - 1, 
		\item NewIterRes = 3 * 1,
\item factorialIter(NewN, Res, NewIterRes) \end{itemize} & 3 > 0. Правда & Прямой ход, решение цели резольвенты\\
	 	\hline
	 	8 & Резольвента:\begin{itemize} \item NewIterRes = 3 * 1,
\item factorialIter(2, Res, NewIterRes) \end{itemize} & NewN = 3 - 1.\linebreak \textbf{Подстановка:} \{N=3, Res=Res, IterRes=1, NewN=2\} & Прямой ход, решение цели резольвенты\\
	 	\hline
	 	9 & Резольвента:\begin{itemize} \item factorialIter(2, Res, 3) \end{itemize} & NewIterRes = 3 * 1.\linebreak \textbf{Подстановка:} \{N=3, Res=Res, IterRes=1, NewN=2, NewIterRes=3\} & Прямой ход, решение цели резольвенты\\
		\hline
		10 & Резольвента:\begin{itemize} \item  2 = 0, \item !, \item Res = 3, \item ! \end{itemize} & factorialIter(2, Res, 3) = factorialIter(N, Res, IterRes). Унификация успешна\linebreak \textbf{Подстановка:} \{N=3, Res=Res, IterRes=1, NewN=2, NewIterRes=3, N=2, Res=Res, IterRes=3\} & Прямой ход, решение цели резольвенты\\
		\hline
	\end{tabular}
\end{table}

\begin{table}[H]
	\begin{tabular}{|p{0.8cm\small}|p{4.7cm\small}|p{5.7cm\small}|p{4cm\small}|}	
		\hline
		№ шага & Состояния резольвенты и вывод: дальнейшие действия (почему?) & Для каких термов запускается алгоритм унификации: T1=T2 и каков результат (подстановка) & Дальнейшие действия: прямой ход или откат (почему и к чему приводит?)\\
		\hline
		11 & Резольвента:\begin{itemize}\item !, \item Res = 3, \item ! \end{itemize} & 2 = 0. Ложь & Откат, переход к следующему шагу относительно 10\\
		\hline
		12 & Резольвента:\begin{itemize} \item 2 > 0, \item NewN = 2 - 1, 
		\item NewIterRes = 2 * 3,
\item factorialIter(NewN, Res, NewIterRes)\end{itemize} & factorialIter(2, Res, 3) = factorialIter(N, Res, IterRes). Унификация успешна\linebreak \textbf{Подстановка:} \{N=3, Res=Res, IterRes=1, NewN=2, NewIterRes=3, N=2, Res=Res, IterRes=3\} & Прямой ход, решение цели резольвенты\\
		\hline
		13 & Резольвента:\begin{itemize} \item NewN = 2 - 1, 
			\item NewIterRes = 2 * 3,
\item factorialIter(NewN, Res, NewIterRes)\end{itemize} & 2 > 0. Правда & Прямой ход, решение цели резольвенты\\
		\hline
		14 & Резольвента:\begin{itemize} \item NewIterRes = 2 * 3,
\item factorialIter(1, Res, NewIterRes)\end{itemize} & NewN = 2 - 1.\linebreak \textbf{Подстановка:} \{N=3, Res=Res, IterRes=1, NewN=2, NewIterRes=3, N=2, Res=Res, IterRes=3, NewN=1\} & Прямой ход, решение цели резольвенты\\
		\hline
		15 & Резольвента:\begin{itemize} \item factorialIter(1, Res, 6)\end{itemize} & NewIterRes = 2 * 3.\linebreak \textbf{Подстановка:} \{N=3, Res=Res, IterRes=1, NewN=2, NewIterRes=3, N=2, Res=Res, IterRes=3, NewN=1, NewIterRes=6\} & Прямой ход, решение цели резольвенты\\
		\hline
	\end{tabular}
\end{table}

\begin{table}[H]
	\begin{tabular}{|p{0.8cm\small}|p{4.7cm\small}|p{5.7cm\small}|p{4cm\small}|}	
		\hline
		№ шага & Состояния резольвенты и вывод: дальнейшие действия (почему?) & Для каких термов запускается алгоритм унификации: T1=T2 и каков результат (подстановка) & Дальнейшие действия: прямой ход или откат (почему и к чему приводит?)\\
		\hline
		16 & Резольвента:\begin{itemize} \item 1 = 0, \item !, \item Res = 6, \item ! \end{itemize} & factorialIter(1, Res, 6) = factorialIter(N, Res, IterRes). Унификация успешна\linebreak \textbf{Подстановка:} \{..., NewN=1, NewIterRes=6, N=1, Res=Res, IterRes=6\} & Прямой ход, решение цели резольвенты\\
		\hline
		17 & Резольвента:\begin{itemize} \item !, \item Res = 6, \item ! \end{itemize} & 1 = 0. Ложь & Откат, переход к следующему предложению относительно 16\\
		\hline
		18 & Резольвента:\begin{itemize} \item 1 > 0, \item NewN = 1 - 1, 
		\item NewIterRes = 1 * 6,
\item factorialIter(NewN, Res, NewIterRes) \end{itemize} & factorialIter(1, Res, 6) = factorialIter(N, Res, IterRes). Унификация успешна\linebreak \textbf{Подстановка:} \{..., NewN=1, NewIterRes=6, N=1, Res=Res, IterRes=6\} & Прямой ход, решение цели резольвенты\\
		\hline
		19 & Резольвента:\begin{itemize} \item NewN = 1 - 1, 
		\item NewIterRes = 1 * 6,
\item factorialIter(NewN, Res, NewIterRes) \end{itemize} & 1 > 0. Правда & Прямой ход, решение цели резольвенты\\
		\hline
		20 & Резольвента:\begin{itemize} \item NewIterRes = 1 * 6,
\item factorialIter(0, Res, NewIterRes) \end{itemize} & NewN = 1 - 1.\linebreak \textbf{Подстановка:} \{..., NewIterRes=6, N=1, Res=Res, IterRes=6, NewN=0\} & Прямой ход, решение цели резольвенты\\
		\hline
	\end{tabular}
\end{table}


\begin{table}[H]
	\begin{tabular}{|p{0.8cm\small}|p{4.7cm\small}|p{5.7cm\small}|p{4cm\small}|}	
		\hline
		№ шага & Состояния резольвенты и вывод: дальнейшие действия (почему?) & Для каких термов запускается алгоритм унификации: T1=T2 и каков результат (подстановка) & Дальнейшие действия: прямой ход или откат (почему и к чему приводит?)\\
		\hline
		21 & Резольвента:\begin{itemize} \item factorialIter(0, Res, 6) \end{itemize} & NewIterRes = 1 * 6.\linebreak \textbf{Подстановка:} \{..., N=1, Res=Res, IterRes=6, NewN=0, NewIterRes=6\} & Прямой ход, решение цели резольвенты\\
		\hline
		22 & Резольвента:\begin{itemize} \item 0 = 0, \item !, \item Res = 6, \item ! \end{itemize} & factorialIter(0, Res, 6) = factorialIter(N, Res, IterRes). Унификация успешна\linebreak \textbf{Подстановка:} \{..., NewN=1, NewIterRes=6, N=1, Res=Res, IterRes=6, N=0, Res=Res, IterRes=6\} & Прямой ход, решение цели резольвенты\\
		\hline
		23 & Резольвента:\begin{itemize} \item !, \item Res = 6 ,\item \end{itemize} & 0 = 0. Правда & Прямой ход, решение цели резольвенты\\
		\hline
		24 & Резольвента:\begin{itemize} \item Res = 6, \item ! \end{itemize} & !. Отсечение 22, 23 & Прямой ход, решение цели резольвенты\\
		\hline
		25 & Резольвента:\begin{itemize} \item ! \end{itemize} & Res = 6.\linebreak \textbf{Подстановка:} \{..., N=0, Res=6, IterRes=6\} & Прямой ход, решение цели резольвенты\\
		\hline
		26 & Резольвента: пуста\linebreak \textbf{Вывод:} Res=6 & !. Отсечение 22-25 & Откат к пункту 22, завершение работы \\
		\hline
	\end{tabular}
\end{table}

Запрос для задания fibon(3, Res).
\begin{table}[H]
	\begin{tabular}{|p{0.8cm\small}|p{4.7cm\small}|p{5.7cm\small}|p{4cm\small}|}	
		\hline
		№ шага & Состояния резольвенты и вывод: дальнейшие действия (почему?) & Для каких термов запускается алгоритм унификации: T1=T2 и каков результат (подстановка) & Дальнейшие действия: прямой ход или откат (почему и к чему приводит?)\\
		\hline
		1 & Резольвента:\begin{itemize} \item fibon(3, Res) \end{itemize} & fibon(3, Res) = factorialIter(N, Res, IterRes). Унификация неуспешна & Прямой ход, переход к следующему предложению \\
		\hline
		2-6 & ... & ... & ...\\
		\hline
	\end{tabular}
\end{table}

\begin{table}[H]
	\begin{tabular}{|p{0.8cm\small}|p{4.7cm\small}|p{5.7cm\small}|p{4cm\small}|}	
		\hline
		№ шага & Состояния резольвенты и вывод: дальнейшие действия (почему?) & Для каких термов запускается алгоритм унификации: T1=T2 и каков результат (подстановка) & Дальнейшие действия: прямой ход или откат (почему и к чему приводит?)\\
		\hline
		7 & Резольвента:\begin{itemize} \item fibonIter(3, Res, 0, 1) \end{itemize} & fibon(3, Res) = fibon(N, Res).Унификация успешна\linebreak \textbf{Подстановка:} \{N=3, Res=Res\} & Прямой ход, решение цели резольвенты\\
		\hline
		8 & Резольвента:\begin{itemize} \item fibonIter(3, Res, 0, 1) \end{itemize} & fibonIter(3, Res, 0, 1) = factorialIter(N, Res, IterRes). Унификация неуспешна & Прямой ход, переход к следующему предложению \\
		\hline
		9-10 & ... & ... & ...\\
		\hline
		11 & Резольвента:\begin{itemize} \item 3 = 0, \item !, \item Res = 0, \item ! \end{itemize} & fibonIter(3, Res, 0, 1) = fibonIter(N, Res, FNum, \_).Унификация успешна\linebreak \textbf{Подстановка:} \{N=3, Res=Res, N=3, Res=Res, FNum=0\} & Прямой ход, решение цели резольвенты\\
		\hline
		12 & Резольвента:\begin{itemize} \item !, \item Res = 0, \item ! \end{itemize} & 3 = 0. Ложь & Откат, переход к следующему предложению относительно 11\\
		\hline
		13 & Резольвента:\begin{itemize} \item 3 = 1, \item !, \item Res = 1, \item ! \end{itemize} & fibonIter(3, Res, 0, 1) = fibonIter(N, Res, \_, SNum).Унификация успешна\linebreak \textbf{Подстановка:} \{N=3, Res=Res, N=3, Res=Res, SNum=1\} & Прямой ход, решение цели резольвенты\\
		\hline
		14 & Резольвента:\begin{itemize} \item !, \item Res = 1, \item ! \end{itemize} & 3 = 1. Ложь & Откат, переход к следующему предложению относительно 13\\
		\hline
	\end{tabular}
\end{table}

\begin{table}[H]
	\begin{tabular}{|p{0.8cm\small}|p{4.7cm\small}|p{5.7cm\small}|p{4cm\small}|}	
		\hline
		№ шага & Состояния резольвенты и вывод: дальнейшие действия (почему?) & Для каких термов запускается алгоритм унификации: T1=T2 и каков результат (подстановка) & Дальнейшие действия: прямой ход или откат (почему и к чему приводит?)\\
		\hline
		15 & Резольвента:\begin{itemize} \item 3 > 1, \item NewN = 3 - 1, \item NewFNum = 1, \item NewSNUm = 0 + 1, \item fibonIter(NewN, Res, NewFNum, NewSNum) \end{itemize} & fibonIter(3, Res, 0, 1) = fibonIter(N, Res, FNum, SNum).Унификация успешна\linebreak \textbf{Подстановка:} \{N=3, Res=Res, N=3, Res=Res, FNum=0, SNum=1\} & Прямой ход, решение цели резольвенты\\
		\hline
		16 & Резольвента:\begin{itemize} \item NewN = 3 - 1, \item NewFNum = 1, \item NewSNUm = 0 + 1, \item fibonIter(NewN, Res, NewFNum, NewSNum) \end{itemize} & 3 > 1. Правда & Прямой ход, решение цели резольвенты\\
		\hline
		17 & Резольвента:\begin{itemize} \item NewFNum = 1, \item NewSNUm = 0 + 1, \item fibonIter(2, Res, NewFNum, NewSNum) \end{itemize} & NewN = 3 - 1\linebreak \textbf{Подстановка:} \{N=3, Res=Res, N=3, Res=Res, FNum=0, SNum=1, NewN=2\} & Прямой ход, решение цели резольвенты\\
		\hline
		18 & Резольвента:\begin{itemize} \item NewSNUm = 0 + 1, \item fibonIter(2, Res, 1, NewSNum) \end{itemize} & NewFNum = 1\linebreak \textbf{Подстановка:} \{N=3, Res=Res, N=3, Res=Res, FNum=0, SNum=1, NewN=2, NewFNum=1\} & Прямой ход, решение цели резольвенты\\
		\hline
		19 & Резольвента:\begin{itemize} \item fibonIter(2, Res, 1, 1) \end{itemize} & NewSNum = 0 + 1\linebreak \textbf{Подстановка:} \{N=3, Res=Res, N=3, Res=Res, FNum=0, SNum=1, NewN=2, NewFNum=1, NewSNum=1\} & Прямой ход, решение цели резольвенты\\
		\hline
	\end{tabular}
\end{table}

\begin{table}[H]
	\begin{tabular}{|p{0.8cm\small}|p{4.7cm\small}|p{5.7cm\small}|p{4cm\small}|}	
		\hline
		№ шага & Состояния резольвенты и вывод: дальнейшие действия (почему?) & Для каких термов запускается алгоритм унификации: T1=T2 и каков результат (подстановка) & Дальнейшие действия: прямой ход или откат (почему и к чему приводит?)\\
		\hline
		20-26 & Резольвента:\begin{itemize} \item fibonIter(2, Res, 1, 1) \end{itemize} & Аналогично пунктам 8-14 & \\
		\hline
		27 & Резольвента:\begin{itemize} \item 2 > 1, \item NewN = 2 - 1, \item NewFNum = 1, \item NewSNUm = 1 + 1, \item fibonIter(NewN, Res, NewFNum, NewSNum)\end{itemize}  & fibonIter(2, Res, 1, 1) = fibonIter(N, Res, FNum, SNum). Унификация успешна\linebreak \textbf{Подстановка:} \{N=3, Res=Res, N=3, Res=Res, FNum=0, SNum=1, NewN=2, NewFNum=1, NewSNum=1, N=2, Res=Res, FNum=1, SNum=1\} & Прямой ход, решение цели резольвенты\\
		\hline
		28 & Резольвента:\begin{itemize} \item NewN = 2 - 1, \item NewFNum = 1, \item NewSNUm = 1 + 1, \item fibonIter(NewN, Res, NewFNum, NewSNum)\end{itemize} & 2 > 1. Правда & Прямой ход, решение цели резольвенты\\
		\hline
		29 & Резольвента:\begin{itemize} \item NewFNum = 1, \item NewSNUm = 1 + 1, \item fibonIter(1, Res, NewFNum, NewSNum)\end{itemize} & NewN = 2 - 1\linebreak \textbf{Подстановка:} \{N=3, Res=Res, N=3, Res=Res, FNum=0, SNum=1, NewN=2, NewFNum=1, NewSNum=1, N=2, Res=Res, FNum=1, SNum=1, NewN=1\} & Прямой ход, решение цели резольвенты\\
		\hline
		30 & Резольвента:\begin{itemize} \item NewSNUm = 1 + 1, \item fibonIter(1, Res, 1, NewSNum)\end{itemize} & NewFNum = 1\linebreak \textbf{Подстановка:} \{..., N=2, Res=Res, FNum=1, SNum=1, NewN=1, NewFNum=1\} & Прямой ход, решение цели резольвенты\\
		\hline
	\end{tabular}
\end{table}

\begin{table}[H]
	\begin{tabular}{|p{0.8cm\small}|p{4.7cm\small}|p{5.7cm\small}|p{4cm\small}|}	
		\hline
		№ шага & Состояния резольвенты и вывод: дальнейшие действия (почему?) & Для каких термов запускается алгоритм унификации: T1=T2 и каков результат (подстановка) & Дальнейшие действия: прямой ход или откат (почему и к чему приводит?)\\
		\hline
		31 & Резольвента:\begin{itemize} \item fibonIter(1, Res, 1, 2)\end{itemize} & NewSNUm = 1 + 1\linebreak \textbf{Подстановка:} \{..., N=2, Res=Res, FNum=1, SNum=1, NewN=1, NewFNum=1, NewSNum=2\} & Прямой ход, решение цели резольвенты\\
		\hline
		32-36 & Резольвента:\begin{itemize} \item fibonIter(1, Res, 1, 2)\end{itemize} & Аналогично пунктам 8-12 & \\
		\hline
		37 & Резольвента:\begin{itemize} \item 1 = 1, \item !, \item Res = 2, \item !\end{itemize} & fibonIter(1, Res, 1, 2) = fibonIter(N, Res, \_, SNum). Унификация успешна\linebreak \textbf{Подстановка:} \{..., N=2, Res=Res, FNum=1, SNum=1, NewN=1, NewFNum=1, NewSNum=2, N=1, Res=Res, SNum=2\} & Прямой ход, решение цели резольвенты\\
		\hline
		38 & Резольвента:\begin{itemize} \item !, \item Res = 2, \item !\end{itemize} & 1 = 1. Правда & Прямой ход, решение цели резольвенты\\
		\hline
		39 & Резольвента:\begin{itemize} \item Res = 2, \item !\end{itemize} & !. Отсечение 37, 38 & Прямой ход, решение цели резольвенты\\
		\hline
		40 & Резольвента:\begin{itemize} \item !\end{itemize} & Res = 2\linebreak \textbf{Подстановка:} \{..., N=2, Res=2, FNum=1, SNum=1, NewN=1, NewFNum=1, NewSNum=2\} & Прямой ход, решение цели резольвенты\\
		\hline
		41 & Резольвента: пуста\linebreak \textbf{Вывод:} Res=2 & !. Отсечение 37-40 & Откат к пункту 37, завершение работы\\
		\hline
	\end{tabular}
\end{table}
		
\end{document}